% Comparación del mejor makespan acumulado de Seteo 4 y Simulated Annealing.
% Los datos provienen de docs/informe/Tarea2/datos/iteraciones_*_pasos.csv,
% que contienen únicamente los puntos donde cambia el mejor makespan.
% Las verticales marcan la primera iteración con makespan 2794:
% 12.867 para Seteo 4 y 14.479 para Simulated Annealing.
% Si se regeneran las trazas con otra semilla, actualizar ambos valores.

\definecolor{agcolor}{HTML}{1F4E79}
\definecolor{sacolor}{HTML}{C55A11}
\definecolor{optcolor}{HTML}{7F7F7F}

\begin{tikzpicture}
\begin{axis}[
    width=15cm,
    height=8.2cm,
    xlabel={Iteración (evaluaciones acumuladas)},
    ylabel={Mejor makespan acumulado},
    xmin=0,
    xmax=27550,
    ymin=2750,
    ymax=4850,
    xtick={0,5000,10000,15000,20000,25000},
    ytick={2800,3000,3400,3800,4200,4600},
    scaled x ticks=false,
    scaled y ticks=false,
    xticklabel style={
        font=\scriptsize,
        /pgf/number format/use comma=false,
        /pgf/number format/1000 sep={.},
    },
    yticklabel style={
        font=\scriptsize,
        /pgf/number format/use comma=false,
        /pgf/number format/1000 sep={.},
    },
    grid=major,
    major grid style={draw=black!8, line width=0.35pt},
    axis background/.style={fill=black!1},
    axis line style={draw=black!55},
    tick style={draw=black!55},
    label style={font=\small},
    legend style={
        at={(0.5,1.035)},
        anchor=south,
        legend columns=2,
        column sep=1.6em,
        draw=none,
        fill=none,
        font=\scriptsize,
    },
    legend cell align=left,
]

% Intervalo entre el primer óptimo de cada heurística.
\fill[black!3]
    (axis cs:12867,2794) rectangle (axis cs:14479,4850);

% Instante en que cada heurística alcanza el óptimo 2794.
\draw[agcolor, line width=0.9pt, dash pattern=on 3pt off 2pt]
    (axis cs:12867,2794) -- (axis cs:12867,4850);
\draw[sacolor, line width=0.9pt, dash pattern=on 1.2pt off 1.8pt]
    (axis cs:14479,2794) -- (axis cs:14479,4850);

% Trayectorias del mejor makespan acumulado.
\addplot[agcolor, line width=1.3pt, const plot]
    table[col sep=comma, x=iteracion, y=makespan]
    {Tarea2/datos/iteraciones_dificil_genetico_seteo4_pasos.csv};
\addlegendentry{Seteo 4 (AG memético)}

\addplot[sacolor, line width=1.3pt, dashed, const plot]
    table[col sep=comma, x=iteracion, y=makespan]
    {Tarea2/datos/iteraciones_dificil_sa_27550_pasos.csv};
\addlegendentry{Simulated Annealing}

% Puntos donde cada heurística alcanza por primera vez 2794.
\addplot[only marks, mark=*, mark size=2.4pt, agcolor, forget plot]
    coordinates {(12867,2794)};
\addplot[only marks, mark=square*, mark size=2.1pt, sacolor, forget plot]
    coordinates {(14479,2794)};

% Línea de referencia del óptimo.
\draw[optcolor, densely dotted, line width=0.9pt]
    (axis cs:0,2794) -- (axis cs:27550,2794);
\node[
    font=\scriptsize,
    text=black!65,
    fill=white,
    inner sep=1.5pt,
    anchor=south west,
] at (axis cs:220,2798) {óptimo 2794};

% Etiquetas de los instantes de primer óptimo.
\node[anchor=east, font=\scriptsize, text=agcolor]
    at (axis cs:12690,4600) {\textbf{Seteo 4}: 12.867};
\node[anchor=west, font=\scriptsize, text=sacolor]
    at (axis cs:14660,4300) {\textbf{SA}: 14.479};

\end{axis}
\end{tikzpicture}
